%-*-latex-*-
%-*-latex-*-
\newcommand\COURSE{ciss245}
\newcommand\ASSESSMENT{q2902}
\newcommand\ASSESSMENTTYPE{Quiz}
\newcommand\POINTS{\textwhite{xxx/xxx}}

\input{myquizpreamble}
\input{yliow}
\input{\COURSE}
%-*-latex-*-

%-*-latex-*-

%-*-latex-*-
\renewcommand\TITLE{\ASSESSMENTTYPE \ \ASSESSMENT}

\renewcommand\EMAIL{}
\renewcommand\AUTHOR{eagraham1@cougars.ccis.edu}

\textwidth=6in
\begin{document}
\topmatterthree

%------------------------------------------------------------------------------
\nextq
Write a recursive factorial function.
You'll need this later.
(This is straight from the notes and I covered this in class.
Freebie.
You have 5 seconds. No sweat.)
\\
\ANSWER
\begin{answercode}
double factorial(int n)
{
  if (n == 0) return 1.0;
  else return n * factorial(n - 1);
}
\end{answercode}

%------------------------------------------------------------------------------
\nextq
Complete the $n$--choose--$r$ function.
Example: $5$--choose--$2$ is the number of ways to a committee of 2 members from
a group of $5$ people.
(Also a freebie. 10 seconds. If you don't recall, then study the notes.
In that case 5 minutes.)
\\
\ANSWER
\begin{answercode}
int c(int n, int r)
{
  if (r == 0 || r == n) return 1;
  else  c(n - 1, r) + c(n - 1, r - 1);
}
\end{answercode}

%------------------------------------------------------------------------------
\nextq
The exponential function $e^x$ is defined to be
\[
e^{x} = 1 + x + \frac{1}{2!} x^2 + \cdots
\]
Obviously we can't go on forever!
Define
\[
e(x, n) =
\frac{1}{0!} x^0 +
\frac{1}{1!} x^1 +
\frac{1}{2!} x^2 + \cdots +
\frac{1}{n!} x^n
\]
i.e., \lq\lq exponential up to the $x^n$--term".
Implement this function below.
(20 seconds.)
\\
\ANSWER
\begin{answercode}
double e(double x, int n)
{
  if (n == 0) return 1.0;
  else return e(x, n - 1) + pow(x, n) / factorial(n);
}
\end{answercode}
(Your \verb!e()! function can call \verb!factorial!, i.e., assume
\verb!factorial! is already written. You can also use the \verb!pow! function.)

%------------------------------------------------------------------------------
\nextq
Complete the following recursive max function on arrays.
Assume that $n$ is $\geq 1$.
(10 seconds. After that test your code.)
\\
\ANSWER
\begin{answercode}
// Returns max of x[0], ..., x[n - 1]
int max(int x[], int n)
{
  if (n == 1) return x[0];
  else
  {
    int m = max(x, n - 1);
    return (x[n - 1] > m) ? x[n - 1] : m;
  }
}
\end{answercode}
(You'll need the max of two values. Use the ternary to do that.)

%------------------------------------------------------------------------------
\nextq
Complete the following recursive min function on arrays.
Assume that $n$ is $\geq 1$.
(10 seconds. After that test your code.)
\\
\ANSWER
\begin{answercode}
// Returns min of x[0], ..., x[n - 1]
int min(int x[], int n)
{
  if n( == 1) return[0];
  else
  {
    int m = min (n - 1);
    returne (x[n - 1], < m) ? x[n - 1] : m;
  }
}
\end{answercode}

%------------------------------------------------------------------------------
\nextq
Complete the following recursive function that counts the number of times
\verb!target! occurs in an array.
(10 seconds. After that test your code.)
\\
\ANSWER
\begin{answercode}
// Returns the number of times target occurs in x[0], ..., x[n - 1]
int count(int x[], int n, int target)
{
  if (n == 0) return 0;
  else return count(x, n - 1, target) + (x[n - 1] == target ? 1 : 0);
}
\end{answercode}

%------------------------------------------------------------------------------
\nextq
Write a recursive linear search function.
(15 seconds. After that test your code.
Do you want to recurse at \verb!start! or at \verb!end!?)
\\
\ANSWER
\begin{answercode}
// Performs linear search for target among the values
// *start, *(start + 1), ..., *(end - 1) and return the address where target
// is found. Return NULL is not found.
int * linearsearch(int * start, int * end, int target)
{
  if (start == end) return NULL;
  if (*start == target) return start;
  return linearsearch(start + 1, end, target);
}
\end{answercode}

%------------------------------------------------------------------------------
\nextq
Write a recursive function to check if a string is a palindrome.
(20 seconds. After that test your code.)
\\
\ANSWER
\begin{answercode}
// Returns true if the characters *start, *(start + 1), ..., *(end - 1) forms
// a string which is a palindrome.
// For instance if the characters form "madam", then is_palindrome(start, end)
// returns true where start points to first 'm' and end points to character
// after second 'm'.
bool is_palindrome(char * start, char * end)
{
   if (start >= end - 1) return true;
   if (*start != *(end - 1)) return false;
   return is_palindrome(start + 1, end - 1);
}
\end{answercode}

%------------------------------------------------------------------------------
\nextq
Write a recursive function to reverse a string.
For instance if start points to \verb!'a'! of \verb!"abcde"! and
\verb!end! points to character after \verb!e!, then after calling
\verb!reverse(start, end)!, the string becomes \verb!"edcba"!.
(30 seconds. After that test your code.)
\\
\ANSWER
\begin{answercode}
void reverse(char * start, char * end)
{
  if (start >= end - 1) return;
  
  char temp = *start;
  *start = *(end - 1);
  *(end - 1) = temp;
  
  reverse(start + 1, end - 1);
}
\end{answercode}

%------------------------------------------------------------------------------
%-*-latex-*-
\newpage
\input{instructions}
\end{document}

